\section{Jacobi fields and uniform comparison of exponential Jacobians}
\label{sec:spectral}

This section proves two estimates that remain uniform near conjugate
endpoints: a positive divided-Hessian difference after pullback by an
exponential differential, and a comparison of exponential Jacobians
along an extended affine segment. The radial estimates use only the
fixed compact Riemannian manifold. Weak MTW enters when we compare
different rays on the segment. None of these constants depends on the
scalar perturbations introduced in Section~\ref{sc:section}.
The pullback estimate is Lemma~\ref{spec:pullback}; it supplies the
positive term in the center lower bound of Lemma~\ref{jet:semibound}.
The affine-segment estimate is Lemma~\ref{spec:comparison}.
Corollary~\ref{spec:outward} applies it to a positive barycentric
representation. A fixed positive vertex weight supplies the opposite
extension required for the affine comparison; radial and norm bounds
control its remaining hypotheses. This is the geometric comparison
needed in the final determinant argument.

For a linear map between tangent spaces, singular values and adjoints are
defined using the Riemannian inner products.  We write
\[
  s_1(x,v)\le\cdots\le s_n(x,v)
\]
for the singular values of $(d\exp_x)_v$, including zero singular values
when $v\in G_x$ is conjugate.  Thus
\[
  \sigma_x(v)=\prod_{j=1}^n s_j(x,v).
\]
For a positive definite symmetric matrix $Q$, the notation
$\lambda_j^{\downarrow}(Q)$ lists its eigenvalues in decreasing order.
All matrix identities below are written in parallel orthonormal frames;
their statements are independent of that choice.

\subsection{The divided-action identities}

On a minimizing ray, the divided Hessian $S(T)$ below is $H_T(v)$
from Lemma~\ref{geo:radial-index}. We record its Jacobi form together
with reversal, and distinguish the local action from the global cost
when a branch passes through cut.

\begin{lemma}[Jacobi matrices and reversal]\label{spec:identities}
Let $\gamma(s)=\exp_x(sv)$ be a geodesic and identify its tangent spaces
by parallel transport.  Put $\mathcal R(s)w=R(w,\dot\gamma(s))\dot\gamma(s)$
and let $C,J$ solve
\[
  U''+\mathcal R U=0,
  \qquad C(0)=\Id,\quad C'(0)=0,\quad
  J(0)=0,\quad J'(0)=\Id.
\]
Then
\begin{equation}\label{spec:exp-jacobi}
  (d\exp_x)_{Tv}=\frac{J(T)}{T}
\end{equation}
for $T>0$. Whenever the geodesic to $\gamma(T)$ is nonconjugate,
let $c_{\gamma,T}$ be its smooth local action: one half the squared
length of the geodesic branch obtained by varying the two endpoints.
The source Hessian of its divided action is
\begin{equation}\label{spec:divided-hessian}
  S(T):=D^2_{xx}\frac{c_{\gamma,T}(x,\gamma(T))}{T}
       =J(T)^{-1}C(T),
  \qquad
  S'(T)=-J(T)^{-1}J(T)^{-*}.
\end{equation}
The endpoint is held fixed when taking $D^2_{xx}$; the derivative in $T$
then follows the displayed geodesic.  The identities apply to a smooth
action branch whether or not it minimizes. Before cut, this action is
the cost $c$; at a nonconjugate minimizing cut vector, it is the chosen
smooth upper branch of $c$.

If $y=\exp_xv$ and $\bar v=-P_{x\to y}v$ is the reverse initial velocity,
then
\begin{equation}\label{spec:reversal}
  (d\exp_y)_{\bar v}=\bigl((d\exp_x)_v\bigr)^*,
  \qquad \sigma_y(\bar v)=\sigma_x(v),
\end{equation}
with the indicated source and target tangent spaces.  Finally, for
$v\in I(x)$,
\begin{equation}\label{spec:radial-identity}
  A_x(v)v=v.
\end{equation}
\end{lemma}

\begin{proof}
Varying the initial velocity in the geodesic equation gives
\eqref{spec:exp-jacobi}.  Since $\mathcal R$ is symmetric, the fundamental
matrix
\[
  \Phi(T)=
  \begin{pmatrix}C(T)&J(T)\\ C'(T)&J'(T)\end{pmatrix}
\]
is symplectic.  In particular,
\begin{equation}\label{spec:wronskians}
  J^*J'=J'^*J,
  \qquad C^*J'-C'^*J=\Id,
  \qquad CJ^*=JC^*.
\end{equation}
These identities can also be obtained by differentiating their left
sides and using the Jacobi equation; their values at zero give the
displayed constants.

For a variation of the initial point by $a$ with terminal variation zero,
the variational field is $C(s)a+J(s)b$, where
$b=-J(T)^{-1}C(T)a$.  The initial gradient of the divided action is the
negative initial velocity.  Differentiating it in a normal frame at $x$
therefore gives $S(T)=J(T)^{-1}C(T)$.  The last identity in
\eqref{spec:wronskians} shows that this matrix is symmetric.  Moreover,
the first two identities give
\[
  J^*\bigl(C'-J'J^{-1}C\bigr)
    =J^*C'-J'^*C=-\Id.
\]
Consequently
\[
  \frac{d}{dT}(J^{-1}C)
    =J^{-1}\bigl(C'-J'J^{-1}C\bigr)
    =-J^{-1}J^{-*},
\]
including the order of the two inverse factors in
\eqref{spec:divided-hessian}.

For completeness, symplecticity also gives
\[
  \Phi(T)^{-1}=
  \begin{pmatrix}J'(T)^*&-J(T)^*\\-C'(T)^*&C(T)^*\end{pmatrix}.
\]
A Jacobi field for the reversed geodesic with initial position zero and
initial derivative $a$ has, in the forward parametrization, terminal
data $(0,-a)$.  Multiplication by this inverse matrix shows that its
value at the other endpoint is $J(T)^*a$.  The reversed Jacobi matrix
at time $T$ is thus $J(T)^*$.  Dividing both matrices by $T$ as in
\eqref{spec:exp-jacobi} proves \eqref{spec:reversal}; this argument does
not require $J(T)$ to be invertible.

In the radial direction the Jacobi equation has the explicit solutions
$J(s)v=sv$ and $C(s)v=v$. Thus $S(T)v=v/T$. For $v\in I(x)$,
taking $T=1$ gives \eqref{spec:radial-identity}; the case $v=0$
is immediate.
\end{proof}

\subsection{Radial shortening at a conjugate endpoint}

\begin{lemma}[Uniform radial transversality]\label{spec:radial}
There are constants $c_{\rm r}>0$ and $C_{\rm r}<\infty$, depending only
on $(M,g)$, such that for $x\in M$, $v\in G_x$, $0<\delta\le1/2$, and
$1\le j\le n$,
\begin{equation}\label{spec:singular-shortening}
  c_{\rm r}\bigl(s_j(x,v)+\delta\bigr)
   \le s_j(x,(1-\delta)v)
   \le C_{\rm r}\bigl(s_j(x,v)+\delta\bigr).
\end{equation}
There is also a positive definite form $Q_{x,v,\delta}$ on $T_xM$ such
that, on putting
\[
  N_{x,\delta}(v)
       =-\delta\bigl(A_x((1-\delta)v)-\Id\bigr),
\]
one has
\begin{align}
  &\|N_{x,\delta}(v)-Q_{x,v,\delta}\|\le C_{\rm r}\delta,
       \label{spec:bounded-background}\\
  &c_{\rm r}\frac{\delta}{s_j(x,v)+\delta}
     \le\lambda_j^{\downarrow}(Q_{x,v,\delta})
     \le C_{\rm r}\frac{\delta}{s_j(x,v)+\delta}.
       \label{spec:hessian-spectrum}
\end{align}
In particular, uniformly on this set,
\begin{equation}\label{spec:bounded-N}
  \|N_{x,\delta}(v)\|\le C_{\rm r},\qquad
  N_{x,\delta}(v)\succeq-C_{\rm r}\delta\Id,
  \qquad N_{x,\delta}(v)v=0.
\end{equation}
No condition on the multiplicity of conjugacy is required.
\end{lemma}

\begin{proof}
The bundle
\[
  \mathcal G=\{(x,v):x\in M,\ v\in G_x\}
\]
is compact: its defining distance equality is closed, and $|v|$ is
bounded by $\operatorname{diam}(M)$.  For each such vector consider
$\gamma(s)=\exp_x(sv)$ on $[0,1]$.  Fix the splitting time $h=1/4$ and
allow the final time $T$ to vary in $[1/2,1]$.
The prefix $hv$ belongs to $I(x)$, since it extends by the factor
$1/h=4$ to the original minimizing ray. For the suffix, reverse the
ray at $y_T=\gamma(T)$ and put
\[
 w_T=-(T-h)\dot\gamma(T),\qquad
 a_T=\frac{T}{T-h}\ge\frac43.
\]
Then $\exp_{y_T}w_T=\gamma(h)$ and
$d(y_T,\exp_{y_T}(a_Tw_T))=a_T|w_T|$, so $w_T\in I(y_T)$.
The reversed suffix is therefore before cut; reversal preserves
uniqueness and, by Lemma~\ref{spec:identities}, nonconjugacy.
The parameter family is compact and its image lies in the open
injectivity domain. Thus both piece differentials and their inverses
have uniform bounds. Their time lengths are bounded below by $1/4$.

Let $C,J$ be the matrices from Lemma~\ref{spec:identities}.  Starting
at time $h$, let $C_h(s),J_h(s)$ have the same initial data and solve
the Jacobi equation along $\gamma(h+s)$.  Smooth dependence on
$(x,v,T)$ and the preceding nonconjugacy show that
\begin{equation}\label{spec:piece-bounds}
  J(h),\ J(h)^{-1},\ J_h(T-h),\ J_h(T-h)^{-1},\ C(h),\ J'(h)
\end{equation}
all have uniformly bounded norms.  This compactness assertion also
includes $v=0$, for which $J(s)=s\Id$.

Write $z=\gamma(h)$ and $y_T=\gamma(T)$.  The Hessian in the joining
point of the broken action
\[
  z'\longmapsto \frac{c(x,z')}{h}
                     +\frac{c(z',y_T)}{T-h}
\]
at $z$ is
\begin{equation}\label{spec:middle-Hessian}
  D(T)=J'(h)J(h)^{-1}
            +J_h(T-h)^{-1}C_h(T-h).
\end{equation}
The first term is the terminal Hessian of the first piece; the second
is the source Hessian of the second piece.  At $T=1$, the broken
action has a minimum at $z$.  To see this directly, the triangle
inequality and Cauchy--Schwarz give
\[
  \frac{d(x,z')^2}{2h}+\frac{d(z',y_1)^2}{2(1-h)}
    \ge\frac{(d(x,z')+d(z',y_1))^2}{2}
    \ge\frac{d(x,y_1)^2}{2},
\]
with equality at $z$.  Hence $D(1)\succeq0$.

Only the second term in \eqref{spec:middle-Hessian} depends on $T$.
Applying \eqref{spec:divided-hessian} with initial point $z$ yields
\begin{equation}\label{spec:middle-crossing}
\begin{split}
  D(1-\delta)&=D(1)+\delta G_\delta,\\
  G_\delta&=\frac1\delta\int_{1-\delta}^{1}
       J_h(s-h)^{-1}J_h(s-h)^{-*}\,ds,
       \qquad c\Id\preceq G_\delta\preceq C\Id.
\end{split}
\end{equation}
The bounds on $G_\delta$ are precisely the bounds in
\eqref{spec:piece-bounds}: the smallest eigenvalue of the integrand
is $\|J_h(s-h)\|^{-2}$, and its largest eigenvalue is
$\|J_h(s-h)^{-1}\|^2$.  Thus the crossing is positive on the whole
space, including any null space of $D(1)$.

Jacobi propagation through the joining point gives the exact
factorization
\begin{equation}\label{spec:Jacobi-factorization}
\begin{split}
  J(T)
    &=C_h(T-h)J(h)+J_h(T-h)J'(h)\\
    &=J_h(T-h)D(T)J(h).
\end{split}
\end{equation}
This formula is also valid at $T=1$ when either of the two full
matrices is singular.  The splitting point is kept fixed while $T$
varies; differentiating a moving midpoint is unnecessary.  Taking
$h=1/2$ gives the usual midpoint factorization on $T\in[3/4,1]$.
The choice $h=1/4$ covers all the parameters in the statement with
one set of bounds.

Write $\lambda_1(D)\le\cdots\le\lambda_n(D)$ for the increasing
eigenvalues of a positive semidefinite matrix $D$.  The min--max
characterization applied to \eqref{spec:middle-crossing} gives
\[
  \lambda_j(D(1))+c\delta
     \le\lambda_j(D(1-\delta))
     \le\lambda_j(D(1))+C\delta.
\]
Multiplication of a matrix on either side by a uniformly bounded
invertible matrix changes each ordered singular value by at most a
uniform multiplicative factor.  Apply this fact to
\eqref{spec:Jacobi-factorization}, use
$(d\exp_x)_{Tv}=J(T)/T$, and note $1/2\le T\le1$.
At $T=1$, the singular values of $J(1)$ are comparable to the
eigenvalues of $D(1)$; at $T=1-\delta$, they are comparable to the
eigenvalues of $D(1-\delta)$.  These observations prove
\eqref{spec:singular-shortening} with uniform constants.

We give the Hessian calculation as well, because it identifies the
bounded term that must not be discarded at a radial eigenvector.
Set
\[
  P=J(h)^{-1}C(h),\qquad E=J(h)^{-*}.
\]
The source--join mixed Hessian of the first divided action is
$-J(h)^{-1}=-E^*$.  For $T<1$, eliminating the joining-point
variation in the Hessian of the broken action gives
\begin{equation}\label{spec:Schur-action}
  \frac{A_x(Tv)}{T}=P-E^*D(T)^{-1}E.
\end{equation}
To verify the elimination, write the quadratic form for source
variation $a$ and joining variation $b$ as
$a^*Pa-2a^*E^*b+b^*D(T)b$.  Its unique stationary value in $b$ is
at $b=D(T)^{-1}Ea$ and has value
$a^*(P-E^*D(T)^{-1}E)a$.  This is the Hessian of the stationary
composed action.  Positivity of $D(T)$ follows from
\eqref{spec:middle-crossing}; the composed action is the smooth
cost branch of the shortened minimizing ray.

For $T=1-\delta$, formula \eqref{spec:Schur-action} gives
\begin{equation}\label{spec:Q-definition}
\begin{split}
  N_{x,\delta}(v)&=Q_{x,v,\delta}
                         +\delta(\Id-TP),\\
  Q_{x,v,\delta}&=T\delta E^*D(T)^{-1}E.
\end{split}
\end{equation}
The second form is positive definite.  The bounds on $P,E,E^{-1}$,
the min--max inequalities just proved, and inversion of $D(T)$ give
\eqref{spec:bounded-background} and \eqref{spec:hessian-spectrum}.
They imply the norm and lower bounds in \eqref{spec:bounded-N}.
The remaining identity follows from
\eqref{spec:radial-identity}.  In particular, for $v\ne0$ the radial
eigenvalue of $N_{x,\delta}(v)$ is exactly zero, while the radial
singular value of $(d\exp_x)_v$ is one.  These two facts are consistent
because the comparison with $Q_{x,v,\delta}$ includes the bounded
background $\delta(\Id-TP)$.
\end{proof}

\begin{corollary}[Radial comparison of exponential determinants]
\label{spec:radial-determinants}
With a constant depending only on $(M,g)$, if $v\in G_x$ and
$1/2\le\theta\le1$, then
\begin{equation}\label{spec:radial-det}
  \sigma_x(\theta v)\ge c\sigma_x(v).
\end{equation}
Suppose $y=\exp_xp$, $p\in G_x$, and
$\bar p=-P_{x\to y}p$.  If $l\in[1/2,1)$ and $\bar p/l\in G_y$,
then
\begin{equation}\label{spec:center-true-det}
  \sigma_y(\bar p/l)\le C\sigma_x(p).
\end{equation}
If $l_+\in[1,2]$, then
\begin{equation}\label{spec:outer-true-det}
  \sigma_y(\bar p/l_+)\ge c\sigma_x(p).
\end{equation}
Comparison~\eqref{spec:center-true-det} permits the longer reversed
vector $\bar p/l$ to have a conjugate endpoint.
\end{corollary}

\begin{proof}
For $\theta<1$, multiply the lower bounds in
\eqref{spec:singular-shortening} with $\delta=1-\theta$ and discard
the nonnegative terms $\delta$.  The case $\theta=1$ is immediate.
For \eqref{spec:center-true-det}, apply \eqref{spec:radial-det} to
the longer vector $\bar p/l$ with shortening factor $l$, and then
use \eqref{spec:reversal}.  For \eqref{spec:outer-true-det}, apply
it to $\bar p$ with shortening factor $1/l_+$ and use the same
reversal identity.  Thus the direction of the two inequalities
reflects, respectively, a longer and a shorter true reversed ray.
\end{proof}

\subsection{A uniform gain from an extra minimizing extension}

\begin{lemma}[Isotropic gain after pullback]\label{spec:pullback}
There is $c_{\rm p}>0$, depending only on $(M,g)$, with the following
property.  Let $y\in M$, $r\in T_yM$, and $1<T_*\le2$ satisfy
$T_*r\in G_y$.  Put
\[
  \gamma(s)=\exp_y(sr),\quad x=\gamma(1),\quad
  p=-P_{y\to x}r,\quad L=(d\exp_x)_p,
  \qquad \tau=1+\frac{T_*-1}{2}.
\]
Then all costs appearing in
\begin{equation}\label{spec:half-extension-gain}
  L^*\left[A_y(r)-\frac1\tau A_y(\tau r)\right]L
       \succeq c_{\rm p}(T_*-1)\Id
\end{equation}
are smooth.  If the terminal ray at time $T_*$ is nonconjugate,
the same inequality holds with $\tau$ replaced by $T_*$ and its
chosen smooth branch.  In particular, for $T_*=1/l$ this latter
inequality reads
\begin{equation}\label{spec:full-extension-gain}
  L^*\left[A_y(r)-l A_y(r/l)\right]L
       \succeq c_{\rm p}(1-l)\Id,
       \qquad 1/2\le l<1.
\end{equation}
Without any uniformity assertion, the source Hessian of divided
action decreases strictly in the positive definite order between
any two nonconjugate times on a minimizing ray.
\end{lemma}

\begin{proof}
Every time less than $T_*$ lies strictly before the cut time of this
ray, so $r$ and $\tau r$ are in $I(y)$, and the reverse vector $p$
is in $I(x)$.  Let $J(s)$ be the Jacobi matrix based at $y$ with
unit initial derivative.  By reversal,
\begin{equation}\label{spec:pullback-transpose}
  L=J(1)^*.
\end{equation}
Apply Lemma~\ref{spec:radial} to the minimizing vector $T_*r$,
with shortening factor $1/T_*$.  Since
$\delta=(T_*-1)/T_*\ge(T_*-1)/2$, its lower bound gives
\begin{equation}\label{spec:J-one-inverse}
  s_1(J(1))=s_1(y,r)\ge c(T_*-1),
  \qquad \|J(1)^{-1}\|\le\frac{C}{T_*-1}.
\end{equation}
There is also a uniform bound on $\|J'(s)\|$ for $0\le s\le T_*$.
Indeed, $|T_*r|\le\operatorname{diam}(M)$, $T_*\le2$, and smooth
Jacobi evolution on this compact set of initial data bounds its
fundamental matrices.  Therefore, for $1\le s\le\tau$,
\begin{equation}\label{spec:noncommuting-ratio}
\begin{split}
  \|J(s)J(1)^{-1}\|
   &\le 1+\|J(s)-J(1)\|\,\|J(1)^{-1}\|\\
   &\le 1+C\frac{s-1}{T_*-1}\le C.
\end{split}
\end{equation}

Write $S(s)=A_y(sr)/s$.  Lemma~\ref{spec:identities} and
\eqref{spec:pullback-transpose} give
\begin{equation}\label{spec:noncommuting-positive}
  L^*[-S'(s)]L
    =J(1)J(s)^{-1}J(s)^{-*}J(1)^*.
\end{equation}
The matrix on the right is $B_sB_s^*$ with
$B_s=J(1)J(s)^{-1}$.  Its smallest eigenvalue equals
$\|B_s^{-1}\|^{-2}=\|J(s)J(1)^{-1}\|^{-2}$ and is thus bounded
below by a uniform positive constant by
\eqref{spec:noncommuting-ratio}.  This uses the displayed order of
the factors and does not commute any Jacobi matrices.  Integration
over $[1,\tau]$, whose length is $(T_*-1)/2$, proves
\eqref{spec:half-extension-gain}.

If the terminal branch is nonconjugate, $S$ is defined through time
$T_*$ and $-S'(s)$ is positive definite on the remaining interval.
Adding its integral over $[\tau,T_*]$ preserves the lower bound.
Since $T_*-1=(1-l)/l\ge1-l$, this gives
\eqref{spec:full-extension-gain}.  For arbitrary two nonconjugate
times on a minimizing ray, integration of
$-S'(s)=J(s)^{-1}J(s)^{-*}\succ0$ proves strict decrease.  When the
true terminal point is conjugate, only the half-extension formula
is used; no undefined terminal Hessian has been differentiated.
\end{proof}

\subsection{From transverse MTW to comparison along an affine segment}

\begin{lemma}[Uniform comparison along an extended segment]
\label{spec:comparison}
Assume weak MTW.  Fix $a_0>0$ and $B<\infty$.  There is
$c_{\rm s}>0$, depending only on $(M,g),a_0,B$, such that the
following holds.  Let $x\in M$, $p,q\in T_xM$, $q\ne0$, satisfy
\begin{equation}\label{spec:segment-hypotheses}
  p+tq\in G_x\quad(-a_0\le t\le1),\qquad
  p\cdot q>0,\qquad
  \frac{|q|^2}{p\cdot q}\le B.
\end{equation}
Then, for every $1\le j\le n$,
\begin{equation}\label{spec:ordered-comparison}
  s_j(x,p)\ge c_{\rm s}s_j(x,p+q).
\end{equation}
In particular,
\begin{equation}\label{spec:product-comparison}
  \sigma_x(p)\ge c_{\rm s}^n\sigma_x(p+q).
\end{equation}
\end{lemma}

\begin{proof}
If there is a configuration in the statement then $B>0$.  Set
\[
  a=\min\{a_0,(2B)^{-1}\}.
\]
We use only the segment $[-a,1]$.  First suppose that all its
points lie in $I(x)$.  Write
\[
  p(t)=p+tq,\qquad \beta=\frac{|q|^2}{p\cdot q},\qquad
  L(t)=\Id+t\frac{qq^*}{p\cdot q}.
\]
The eigenvalues of $L(t)$ are $1$ on $q^\perp$ and $1+t\beta$
on $\R q$.  Hence
\begin{equation}\label{spec:L-bounds}
  \|L(t)\|\le1+B,\qquad \|L(t)^{-1}\|\le2
       \qquad(-a\le t\le1).
\end{equation}
For $0<\delta\le1/2$, put
\begin{equation}\label{spec:F-definition}
  N_\delta(t)=-\delta\bigl(A_x((1-\delta)p(t))-\Id\bigr),
  \qquad F_\delta(t)=L(t)^*N_\delta(t)L(t).
\end{equation}
These matrices are uniformly bounded by Lemma~\ref{spec:radial}
and \eqref{spec:L-bounds}, and
$F_\delta(t)\succeq-C\delta\Id$.

We next prove full matrix convexity of $F_\delta$.  For any fixed
$z\in T_xM$, the vector
\[
  \xi=z-\frac{z\cdot q}{p\cdot q}p
\]
is perpendicular to $q$.  Moreover,
\[
  \xi-L(t)z=-\frac{z\cdot q}{p\cdot q}p(t).
\]
The last vector belongs to the kernel of $N_\delta(t)$ by the
radial identity.  Symmetry consequently gives the exact equality
\begin{equation}\label{spec:transverse-conjugation}
  z^*F_\delta(t)z=\xi^*N_\delta(t)\xi.
\end{equation}
Weak MTW means that
$t\mapsto A_x((1-\delta)p(t))(\xi,\xi)$ is concave whenever
$\xi\perp q$.  The negative sign in the definition of $N_\delta$
therefore makes the right side of
\eqref{spec:transverse-conjugation} convex.  Since this holds for
every $z$, $F_\delta$ is convex in the Loewner order.

This argument remains uniform when $p\cdot q$ is small.  The
fixed test vector $\xi$ itself might be large, but its value
modulo the radial kernel is $L(t)z$.  Equations
\eqref{spec:L-bounds} and \eqref{spec:transverse-conjugation}
bound the scalar convex function by $C|z|^2$, independently of
the size of $\xi$.  For any compact interval
$K\subset(-a,1)$, a real convex function bounded in absolute
value by $C|z|^2$ has slopes on $K$ bounded by
$2C|z|^2/\operatorname{dist}(K,\{-a,1\})$: sandwich each
secant slope between secants to the two endpoints.  Applying
this estimate to basis vectors and their sums and differences
shows by polarization that every matrix entry of $F_\delta$
has a uniform Lipschitz bound on $K$.

To turn this bounded matrix convexity into a spectral comparison,
we argue by contradiction. The spectral profile
$\delta/(s_i+\delta)$ from \eqref{spec:hessian-spectrum}
detects singular values below the radial scale: it stays positive
when $s_i\ll\delta$ and tends to zero when $s_i\gg\delta$.
At a scale separated from the singular values, the rank of a
limiting matrix therefore counts those below the scale.
A failing ordered singular-value ratio
will provide a radial scale between the center and endpoint values.
We choose it to separate asymptotically from every singular value
at each fixed point of a countable dense set of segment parameters.
The resulting positive semidefinite
matrix limit has a uniform gap between zero and its positive
eigenvalues, forcing constant interior rank. Convexity then prevents
rank loss at the endpoint, whereas the failing ratio will force
precisely such a loss.

Suppose that no uniform constant in
\eqref{spec:ordered-comparison} exists.  There are configurations
$(x_k,p_k,q_k)$ satisfying the same $a_0,B$, an index $j$ fixed
after passing to a subsequence, and
\begin{equation}\label{spec:failed-ratio}
  \frac{s_j(x_k,p_k)}{s_j(x_k,p_k+q_k)}\longrightarrow0.
\end{equation}
For now all the segments are in their open injectivity domains,
so the singular values are positive.  Choose an orthonormal
identification of each $T_{x_k}M$ with $\R^n$.  No smooth choice
in $k$ is needed.  Abbreviate
\[
  s_{i,k}(t)=s_i(x_k,p_k+tq_k),\qquad
  R_k=\log\frac{s_{j,k}(1)}{s_{j,k}(0)}\longrightarrow\infty.
\]
Let $\{t_1,t_2,\ldots\}$ be a fixed countable dense subset of
$(-a,1)$, enlarged to include $0$ and $1$.  We now choose one
radial scale that separates every fixed member of this set.
Take $\log\delta_k$ in the middle half of the interval
\[
  [\log s_{j,k}(0),\log s_{j,k}(1)].
\]
Exclude the intervals of radius $R_k^{1/3}$ centered at
\[
  \log s_{i,k}(t_m),\qquad
  1\le i\le n,\quad 1\le m\le\lfloor R_k^{1/3}\rfloor.
\]
Their total length is at most $2nR_k^{2/3}$, while the available
middle interval has length $R_k/2$.  For all sufficiently large
$k$ there is a remaining point, which we choose as
$\log\delta_k$.  This construction ensures
\begin{align}
  &s_{j,k}(0)/\delta_k\longrightarrow0,
       \qquad s_{j,k}(1)/\delta_k\longrightarrow\infty,
       \label{spec:separating-scale}\\
  &\left|\log\frac{s_{i,k}(t_m)}{\delta_k}\right|
       \longrightarrow\infty
       \quad\text{for every fixed }i,m.
       \label{spec:dense-separation}
\end{align}
All singular values on $\mathcal G$ are uniformly bounded above.
Since $\delta_k\le s_{j,k}(1)e^{-R_k/4}$, we also have
$\delta_k\to0$.  In particular these scales are eventually
admissible in Lemma~\ref{spec:radial}.

Form the matrices $F_k(t)$ in
\eqref{spec:F-definition} using the $k$th segment and scale
$\delta_k$.  Their uniform bounds and interior Lipschitz bounds
give, by successive extraction on compact subintervals, a locally
uniform limit $F(t)$ on $(-a,1)$.  Extract also a limit
$E=\lim_kF_k(1)$.  Both $F(t)$ and $E$ are positive semidefinite,
because their approximants are bounded below by
$-C\delta_k\Id$.  The limit $F$ is matrix convex and bounded.

To identify its possible ranks, use
\eqref{spec:bounded-background} to write
\[
  F_k(t)=L_k(t)^*Q_{x_k,p_k+tq_k,\delta_k}L_k(t)
                     +\mathcal E_k(t),
  \qquad \|\mathcal E_k(t)\|\le C\delta_k.
\]
The positive form in this expression has its $i$th decreasing
eigenvalue between uniform multiples of
\begin{equation}\label{spec:rank-spectral-profile}
  \frac{\delta_k}{s_{i,k}(t)+\delta_k},
\end{equation}
by \eqref{spec:hessian-spectrum} and \eqref{spec:L-bounds}.
At a fixed dense point $t_m$, put $\eta_k=\exp(-R_k^{1/3})$.
For all sufficiently large $k$, depending on $m$,
\eqref{spec:dense-separation} says that
$s_{i,k}(t_m)/\delta_k$ is at most $\eta_k$ or at least
$\eta_k^{-1}$. Weyl's eigenvalue inequality therefore gives
one of the two alternatives
\[
 \lambda_i^\downarrow(F_k(t_m))
       \ge \frac{c_0}{1+\eta_k}-C\delta_k,
 \qquad\text{or}\qquad
 -C\delta_k\le\lambda_i^\downarrow(F_k(t_m))
       \le C_0\eta_k+C\delta_k,
\]
with $c_0,C_0>0$ uniform in $i,m,k$.
Since $F_k(t_m)$ already converges, its ordered eigenvalues converge.
Their limits are therefore either zero or at least $c_0$, even if
the alternatives initially alternate. No simultaneous rate at all
dense points is required. Set $c_*=c_0/2$.

The ordered eigenvalues of a symmetric matrix are continuous in
its entries.  Density of the set $\{t_m\}$ and continuity of
$F(t)$ therefore give
\begin{equation}\label{spec:uniform-limit-gap}
  \operatorname{spec}F(t)\subset\{0\}\cup[c_*,C_*]
        \qquad(-a<t<1)
\end{equation}
after decreasing $c_*$ once if necessary.  Indeed, an eigenvalue
strictly between zero and $c_*$ at an interior point would remain
in that interval in a neighborhood and contradict the dense-set
conclusion.  The rank is now locally constant, because a
continuous eigenvalue cannot pass from zero to a positive value
without entering that gap.  It is consequently constant on the
connected interval $(-a,1)$.

The endpoint also has to be treated; constancy of the interior
rank alone would not suffice.  For every fixed vector $z$, the
bounded scalar convex function $z^*F(t)z$ has a finite limit as
$t\uparrow1$.  Polarization yields a matrix limit $F(1-)$.
The gap \eqref{spec:uniform-limit-gap} persists in this limit,
so its rank equals the interior rank.  For fixed $s<t<1$,
matrix convexity of $F_k$ says
\[
  F_k(t)\preceq\frac{1-t}{1-s}F_k(s)
                         +\frac{t-s}{1-s}F_k(1).
\]
First let $k\to\infty$, then let $t\uparrow1$; this gives
\begin{equation}\label{spec:endpoint-order}
  F(1-)\preceq E.
\end{equation}
Positive semidefinite order implies
$\operatorname{rank}E\ge\operatorname{rank}F(1-)$, since the
kernel of $E$ is contained in the kernel of $F(1-)$.

On the other hand, the first inequality in
\eqref{spec:separating-scale} and the ordering of singular values
make the first $j$ eigenvalues in
\eqref{spec:rank-spectral-profile} at $t=0$ bounded away from
zero.  Thus $\operatorname{rank}F(0)\ge j$.  The second inequality
makes all eigenvalues with decreasing index $i\ge j$ tend to
zero at $t=1$, so $\operatorname{rank}E\le j-1$.  This contradicts
the interior rank constancy and \eqref{spec:endpoint-order}.
The uniform ordered comparison is proved for open segments.

Finally, replace a closed segment in \eqref{spec:segment-hypotheses}
by $(1-\eps)(p+tq)$, with $0<\eps<1$.  It lies in $I(x)$ by
radial shortening.  Its extension interval and the ratio
$|q|^2/(p\cdot q)$ are unchanged, so the same constant applies.
Letting $\eps\downarrow0$ and using continuity of the singular
values proves \eqref{spec:ordered-comparison} on $G_x$.
Multiplication over $j$ proves
\eqref{spec:product-comparison}.  Every bound in the contradiction
argument depended only on $(M,g),a_0,B$, which establishes the
stated uniformity.
\end{proof}

\subsection{Jacobian comparison for a barycentric configuration}

\begin{corollary}[Barycentric configurations]\label{spec:outward}
Assume weak MTW and the convexity of $G_x$ established in
Corollary~\ref{geo:convexity}.  Fix $\mu\in(0,1)$ and positive constants
$\kappa,L_0$.  Suppose
\[
  p=\sum_{i=1}^N m_i p_i,\qquad
  p_i\in G_x,\quad m_i>0,\quad\sum_i m_i=1,
\]
and for one index $i$ there are $D>0$, a unit vector $e$, and
$t_i>0$ such that
\begin{equation}\label{spec:outward-data}
  p_i=p+t_i e,\quad h:=p\cdot e>0,\quad
  h^2\ge\kappa D,\quad
  |p_i|^2-|p|^2\le L_0D,\quad m_i\ge\mu.
\end{equation}
Then
\begin{equation}\label{spec:outward-det}
  \sigma_x(p)\ge c\sigma_x(p_i),
\end{equation}
where $c>0$ depends only on $(M,g),\mu,\kappa,L_0$.
\end{corollary}

\begin{proof}
Set $q=t_ie$.  Expansion of the squared norms in
\eqref{spec:outward-data} gives
\[
  2ht_i+t_i^2\le L_0D,
  \qquad
  \frac{|q|^2}{p\cdot q}=\frac{t_i}{h}
          \le\frac{L_0D}{2h^2}\le\frac{L_0}{2\kappa}.
\]
Because $q\ne0$, $m_i<1$.  The remaining barycenter belongs to
$G_x$ by its convexity and is
\[
  \frac1{1-m_i}\sum_{j\ne i}m_jp_j
       =p-\frac{m_i}{1-m_i}q.
\]
The coefficient $m_i/(1-m_i)$ is at least $m_i\ge\mu$.
Convexity therefore puts $p+s q$ in $G_x$ for every
$s\in[-\mu,1]$.  Lemma~\ref{spec:comparison} applies with
$a_0=\mu$ and $B=L_0/(2\kappa)$ and proves the result.
No upper bound on $m_i/(1-m_i)$ is needed: a shorter common
opposite extension is enough.
\end{proof}
